% Generated by `taskgen course examples` from modular.ntt@fermat (seed 2).
% Do not edit: change examples.yaml and regenerate.
\begin{example}[ТЧП Ферма, $m = 257$, $n = 4$]\label{ex:ntt-fermat}
Вычислить циклическую свёртку последовательностей \[a = (8, 1, 1, 8), \qquad b = (4, 7, 4, 9)\] длины $n = 4$ в кольце $\mathbb{Z}/257\mathbb{Z}$, $257 = 2^{8} + 1$, с помощью теоретико-числового преобразования с корнем $\omega = 16$ из единицы степени $4$. Проверить ответ прямым вычислением свёртки.

\par\noindent\textit{Решение.} $257 = 2^{8} + 1$, поэтому $2^{8} \equiv -1$ и $2^{16} \equiv 1$: $\omega = 16 = 2^{4}$ имеет порядок $16 / 4 = 4$, $\omega^{2} \equiv -1$. Прямое преобразование: $A_k = \sum_j a_j \omega^{jk}$, матрица $(\omega^{jk} \bmod m)$ (вычеты $\pm 2^e$):\[W = \begin{pmatrix}1 & 1 & 1 & 1\\1 & 16 & -1 & -16\\1 & -1 & 1 & -1\\1 & -16 & -1 & 16\end{pmatrix}.\]\[\begin{array}{c|l|r|r}k & A_k = \sum_j \omega^{jk} a_j & & A_k \bmod 257 \\ \hline 0 & 8 + 1 + 1 + 8 & 18 & 18 \\ 1 & 8 + 16 \cdot 1 - 1 - 16 \cdot 8 & -105 & 152 \\ 2 & 8 - 1 + 1 - 8 & 0 & 0 \\ 3 & 8 - 16 \cdot 1 - 1 + 16 \cdot 8 & 119 & 119 \\ \end{array}\]\[\begin{array}{c|l|r|r}k & B_k = \sum_j \omega^{jk} b_j & & B_k \bmod 257 \\ \hline 0 & 4 + 7 + 4 + 9 & 24 & 24 \\ 1 & 4 + 16 \cdot 7 - 4 - 16 \cdot 9 & -32 & 225 \\ 2 & 4 - 7 + 4 - 9 & -8 & 249 \\ 3 & 4 - 16 \cdot 7 - 4 + 16 \cdot 9 & 32 & 32 \\ \end{array}\]Поточечное произведение $C_k = A_k B_k \bmod 257$: $18 \cdot 24 = 432 \equiv 175$; $152 \cdot 225 = 34200 \equiv 19$; $0 \cdot 249 = 0 \equiv 0$; $119 \cdot 32 = 3808 \equiv 210$.\par Обратное преобразование: $c_j = n^{-1} \sum_k C_k \omega^{-jk}$, $\omega^{-1} = \omega^{3} \equiv -16$, $n^{-1} = 4^{-1} \equiv 2^{14} \equiv -2^{6} \equiv 193 \pmod{257}$, матрица $(\omega^{-jk} \bmod m)$:\[W^{-} = \begin{pmatrix}1 & 1 & 1 & 1\\1 & -16 & -1 & 16\\1 & -1 & 1 & -1\\1 & 16 & -1 & -16\end{pmatrix}.\]\[\begin{array}{c|l|r|r|r}j & \sum_k \omega^{-jk} C_k & & \bmod 257 & c_j = n^{-1} \cdot (\ldots) \\ \hline 0 & 175 + 19 + 0 + 210 & 404 & 147 & 101 \\ 1 & 175 - 16 \cdot 19 - 0 + 16 \cdot 210 & 3231 & 147 & 101 \\ 2 & 175 - 19 + 0 - 210 & -54 & 203 & 115 \\ 3 & 175 + 16 \cdot 19 - 0 - 16 \cdot 210 & -2881 & 203 & 115 \\ \end{array}\]Проверка: прямое вычисление циклической свёртки над $\mathbb{Z}$, $c_j = \sum_i a_i b_{(j - i) \bmod n}$, даёт $(101, 101, 115, 115)$; все $c_j < 257$, поэтому ТЧП дало \emph{точную} свёртку над $\mathbb{Z}$.\[c = a * b = (101, 101, 115, 115) \pmod{257}.\]
\end{example}
